\section{Dual-first filter and primal admissibility}\label{sec:dualfirst}
We use the notation of the supplied paper~\cite{dualPIE}: the primal factor and weight are $\Theta_\Delta$ and $\mathcal N_\Delta=\diag(\mathcal X_\Delta,-\mathcal Y_\Delta)$; dual quantities are \emph{underlined}; the duality map is $\mathbf D$ and the signature permutation is $\hat J$. For the delay-free ODE specialization, spatial operators become matrices and $L_{[a,b]}$ becomes the temporal signal space $L_2$. The truncation $p_T$ retains the signal on $[0,T]$. All filters have zero initial state.

For the incremental application, use the paper's uncertainty-channel signals for differences:
\begin{equation}\label{eq:incrementalsignals}
 z_\Delta=q^{(1)}-q^{(2)},\qquad
 w_\Delta=S(q^{(1)})-S(q^{(2)}).
\end{equation}
The vector $S$ applies the sigmoid componentwise. Its slope bound $0\le S_i'\le1$ implies $w_\Delta=D(t)z_\Delta$, where $D(t)$ is diagonal with entries in $[0,1]$. No rate bound on $D$ or equilibrium assumption is imposed. More generally, define the causal uncertainty family $\boldsymbol\Delta=\{\Delta_D:z\mapsto D(\cdot)z\}$ for all measurable such $D$. Every pair of sigmoid trajectories lies on a graph in this family.

For each pair choose $0<\rho_j<1$, $a_j>0$, an integer $n_j\ge1$, and $0\le\beta_j<1$. Set
\begin{align}
 R&=\begin{bmatrix}1&1\\-1&0\end{bmatrix},\qquad
 h_j(s)=\left(\frac{\beta_j s+a_j}{s+a_j}\right)^{n_j},\notag\\
 H_j(s)&=\rho_j h_j(s)R,\qquad \eta_j=\frac{(1-\rho_j)^2}{8},\notag\\
 H&=\blkdiag(H_j),\quad E=\blkdiag(\eta_jI_2),\notag\\
 W&=\blkdiag(t_jI_2),\qquad t_j>0.
\end{align}
Define the alternative factor and weight on the \emph{dual} side first:
\begin{equation}\label{eq:dualfirst}
 \boxed{\begin{aligned}
 \underline\Theta_\Delta&=\begin{bmatrix}I&-H\\0&I\end{bmatrix},\\
 \underline{\mathcal N}_\Delta&=\diag(WE^{-1},-W).
 \end{aligned}}
\end{equation}
Equivalently, take $\underline\Psi_\Delta=\underline\Theta_\Delta$ and $\underline{\mathcal V}_\Delta=\underline{\mathcal N}_\Delta$. For fixed filters the dual weight is affine in the positive pair weights $t_j$.

For a causal LTI realization $H=(A,B,C,D)$, its prime dual is $H'=(A^{\T},C^{\T},B^{\T},D^{\T})$. This differs from the generally anticausal Hilbert-space adjoint. With
\begin{equation*}
 \hat J=\begin{bmatrix}0&I\\-I&0\end{bmatrix},\qquad
 \mathbf D(\Theta)=(\mathcal M_{\hat J}^*\Theta'\mathcal M_{\hat J})^{-1},
\end{equation*}
recover the primal factor and weight using~\cite[Definition 7 and Theorem 3]{dualPIE}:
\begin{equation}\label{eq:recoveredprimal}
 \boxed{\begin{aligned}
 \Theta_\Delta&=\mathbf D(\underline\Theta_\Delta)
 =\begin{bmatrix}I&-H'\\0&I\end{bmatrix},\\
 \mathcal N_\Delta&=-\mathbf D(\underline{\mathcal N}_\Delta)
 =\diag(W^{-1},-EW^{-1}).
 \end{aligned}}
\end{equation}
Here $\mathcal M_K$ denotes memoryless multiplication by $K$. Thus $\mathcal X_\Delta=W^{-1}$, $\mathcal Y_\Delta=EW^{-1}$, and $\underline{\mathcal N}_\Delta=\diag(\mathcal Y_\Delta^{-1},-\mathcal X_\Delta^{-1})$. Involutivity gives $\mathbf D(\Theta_\Delta)=\underline\Theta_\Delta$. Both weights have coercive positive and negative diagonal blocks. Moreover,
\begin{align}
 \underline\Theta_\Delta^{-1}&=\begin{bmatrix}I&H\\0&I\end{bmatrix},&
 \Theta_\Delta^{-1}&=\begin{bmatrix}I&H'\\0&I\end{bmatrix},\notag\\
 \Theta_{\Delta,22}&=I,&
 \Theta_{\Delta,21}P+\Theta_{\Delta,22}&=I.
\end{align}
These are bounded causal inverses for every stable $H$. No inverse of $H$ is used. The MATLAB generator defaults to $\beta=0$, allowing strictly proper $H$.

\subsection{Hardness on the primal sigmoid graph}
\begin{prop}[Hard incremental admissibility]\label{prop:alternative}
The recovered primal pair $(\Theta_\Delta,\mathcal M_{\mathcal N_\Delta})$ is an admissible hard factorization in Definition 7 of~\cite{dualPIE} for the family $\boldsymbol\Delta$. In particular, every pair of sigmoid trajectories in~\eqref{eq:incrementalsignals} satisfies, for every $T\ge0$,
\begin{align}\label{eq:dualfirsthard}
 \mathcal I_\Delta(T)
 &: =\left\langle p_T\Theta_\Delta\begin{bmatrix}z_\Delta\\w_\Delta\end{bmatrix},
 p_T\mathcal M_{\mathcal N_\Delta}\Theta_\Delta\begin{bmatrix}z_\Delta\\w_\Delta\end{bmatrix}\right\rangle_{L_2}\notag\\
 &=\sum_j\frac1{t_j}\left(\|z_{\Delta,j}-H_j'w_{\Delta,j}\|_{2,T}^2
 -\eta_j\|w_{\Delta,j}\|_{2,T}^2\right)\notag\\
 &\ge\sum_j\frac{\eta_j}{t_j}\|z_{\Delta,j}\|_{2,T}^2\ge0.
\end{align}
No bounded causal inverse of $H$ is required.
\end{prop}
\begin{proof}
For one pair define the auxiliary norm
\begin{equation*}
 \nu_-(x)=\max\{|x_1|,|x_2|,|x_1-x_2|\}.
\end{equation*}
This symbol is distinct from the paper's signature weight $\mathcal N_\Delta$. Since $R^{\T}x=[x_1-x_2,x_1]^{\T}$, $\nu_-(R^{\T}x)=\nu_-(x)$. Also $\nu_-(Dx)\le\nu_-(x)$: maximizing $|d_1x_1-d_2x_2|$ over $[0,1]^2$ reduces by convexity to $0,|x_1|,|x_2|,|x_1-x_2|$ at the vertices. Finally,
\begin{equation*}
 \|x\|_2/\sqrt2\le\nu_-(x)\le\sqrt2\|x\|_2.
\end{equation*}
Let $\|f\|_{\nu_-,T}^2=\int_0^T\nu_-(f(t))^2\,dt$. Drop the pair index temporarily. The binomial expansion
\begin{equation*}
 h(s)=\sum_{k=0}^{n}\binom nk\beta^{n-k}(1-\beta)^k
 \left(\frac{a}{s+a}\right)^k
\end{equation*}
has nonnegative coefficients summing to one. Each $k\ge1$ term has a nonnegative unit-integral Erlang kernel; $k=0$ is the identity. Young's convolution inequality, after zero-extending the input outside $[0,T]$, gives $\|hf\|_{\nu_-,T}\le\|f\|_{\nu_-,T}$. This includes $\beta=0$.

Since $H'=\rho h R^{\T}$ and $w_\Delta=D z_\Delta$,
\begin{equation*}
 \|H'w_\Delta\|_{\nu_-,T}\le\rho\|w_\Delta\|_{\nu_-,T}
 \le\rho\|z_\Delta\|_{\nu_-,T}.
\end{equation*}
With $r=z_\Delta-H'w_\Delta$, the reverse triangle inequality gives
\begin{equation*}
 (1-\rho)\|z_\Delta\|_{\nu_-,T}\le\|r\|_{\nu_-,T}.
\end{equation*}
By norm equivalence and the diagonal slope bound,
\begin{align*}
 \|r\|_{2,T}^2&\ge\frac{(1-\rho)^2}{4}\|z_\Delta\|_{2,T}^2
 =2\eta\|z_\Delta\|_{2,T}^2,\\
 \|w_\Delta\|_{2,T}&\le\|z_\Delta\|_{2,T}.
\end{align*}
Subtracting $\eta\|w_\Delta\|_{2,T}^2$ and summing with positive weights $t_j^{-1}$ proves~\eqref{eq:dualfirsthard}. To complete Definition 7, take $\Psi_\Delta=\Theta_\Delta$, $\mathcal V_\Delta=\mathcal N_\Delta$, and $\mathcal U_\Delta=I$: the factorization equality is immediate. The blocks $\mathcal X_\Delta,\mathcal Y_\Delta$ are coercive and the displayed triangular inverses are stable and causal. This establishes all uncertainty-channel hard-factorization requirements.
\end{proof}
Thus the candidate meets the uncertainty-channel conditions of R.2(i) of~\cite{dualPIE}. The remaining plant well-posedness, performance positivity and strict dual plant inequality are separate requirements. If a performance channel is added, its factor and the full input-coordinate map must also satisfy the paper's assumptions.

\subsection{Dual analysis and implementation}
The complementary test is performed on the graph of the LTI dual $P'$:
\begin{align*}
 \underline{\mathcal I}(T)&=
 \left\langle p_T\underline\Theta_\Delta\begin{bmatrix}P'\underline w\\\underline w\end{bmatrix},
 p_T\mathcal M_{\underline{\mathcal N}_\Delta}\underline\Theta_\Delta\begin{bmatrix}P'\underline w\\\underline w\end{bmatrix}
 \right\rangle_{L_2}\\
 &\le-\underline\epsilon\|p_T\underline w\|_{L_2}^2.
\end{align*}
On the illustrative matching plant $P'=H$, $P=H'$, the first filtered output vanishes identically. Hence $\underline{\mathcal I}=-\|\underline w\|_{W,T}^2$, while the primal plant supply is $\mathcal I=-\|w\|_{EW^{-1},T}^2$. Both inequalities are strict, illustrating Theorem 3. This is a test plant, not the PNAS network. For a single unit-weight pair, a strict bound on both graph signals as used by~\cite{suKhong} follows with $\varepsilon_{\mathrm{IO}}=\eta/(1+\rho^2\|R\|_2^2)$.
\begin{lstlisting}[language=Matlab]
addpath('Notes/sigmoid_iiqc')
iqc = sigmoid_dual_iiqc(0.95,100,10); % beta=0
Theta_Delta = iqc.Theta_Delta;
N_Delta = iqc.N_Delta;
Theta_Delta_underlined = iqc.Theta_Delta_underlined;
N_Delta_underlined = iqc.N_Delta_underlined;
iqc4 = sigmoid_dual_iiqc([0.9 0.95],100,10,0,[1 2]);
verify_sigmoid_dual_kyp
\end{lstlisting}
The supplied weights parameterize the dual side; their reciprocals occur in the recovered primal supply. Arbitrary sums of dual multipliers cannot be recovered by simply summing the corresponding primal multipliers, since $\mathbf D$ contains inversion. A general summed dual candidate requires its own admissible hard factorization. The positive pair-weight family above avoids this issue.

\subsection{KYP feasibility test of the alternative multiplier}
For the delay-free ODE case, the paper's dual filtered graph is
\begin{equation*}
 \underline\Theta_\Delta\begin{bmatrix}P'\\I\end{bmatrix}:
 \quad\begin{cases}
 \dot{\underline x}=\underline{\mathcal A}\underline x+\underline{\mathcal B}\underline w,\\
 \underline v=\underline{\mathcal C}\underline x+\underline{\mathcal D}\underline w.
 \end{cases}
\end{equation*}
The ODE specialization of the dual KYP LPI in~\cite[Section VII.B]{dualPIE} seeks a storage matrix $\mathcal P\succeq0$ satisfying
\begin{equation}\label{eq:alternativekyp}
\begin{aligned}
 &\begin{bmatrix}
 \underline{\mathcal A}^{\T}\mathcal P+\mathcal P\underline{\mathcal A}
 &\mathcal P\underline{\mathcal B}\\
 \underline{\mathcal B}^{\T}\mathcal P&\underline\epsilon I
 \end{bmatrix}\\
 &\quad+\begin{bmatrix}\underline{\mathcal C}^{\T}\\\underline{\mathcal D}^{\T}\end{bmatrix}
 \underline{\mathcal N}_\Delta
 \begin{bmatrix}\underline{\mathcal C}&\underline{\mathcal D}\end{bmatrix}
 \preceq0.
\end{aligned}
\end{equation}
Here $\mathcal P$ denotes storage, distinct from the plant $P$. For $V=\underline x^{\T}\mathcal P\underline x$, this gives
\begin{equation*}
 \dot V+\underline v^{\T}\underline{\mathcal N}_\Delta\underline v
 +\underline\epsilon\|\underline w\|^2\le0.
\end{equation*}
Integration from zero initial state and $V(T)\ge0$ proves the strict finite-horizon complementary inequality. Primal sigmoid admissibility is supplied separately by Proposition~\ref{prop:alternative}.

The reusable function \texttt{kyp\_iqc\_graph.m} assembles the full augmented plant/filter realization and solves this SDP using YALMIP and SeDuMi. It accepts any compatible fixed numeric ODE plant. It returns the storage, KYP matrix, solver status, and recomputed eigenvalues. No frequency sampling or state cancellation is used. For numerical conditioning, the inequality is divided by its prescribed positive input margin before solving, then converted back to original units.

The script \texttt{verify\_sigmoid\_dual\_kyp.m} is now a self-contained 38-line example. It constructs the scalar cascade, the paired $H$, the dual factor, and the filtered dual graph directly with MATLAB state-space objects, then solves~\eqref{eq:alternativekyp}. No multiplier generator or helper is needed. The script defaults to $\rho=0.95$, $a=1$, $n=10$, $\beta=0$, and $\underline\epsilon=0.5$. Change $\beta$ to $0.5$ to test the biproper case, or replace \texttt{Pprime = H} with another fixed dual plant. The storage and KYP eigenvalues are checked after solving. Control System Toolbox, YALMIP and SeDuMi are required.

Earlier comparison runs found feasible primal and dual KYP certificates for both $\beta$ values, while the optimized diagonal static sector KYP tests were infeasible. Those results remain in \texttt{dual\_kyp\_results.mat} and \texttt{dual\_kyp\_run.log}; the simplified example focuses on one dual KYP inequality.

\subsection{Additional trajectory and algebra tests}
The script \texttt{verify\_sigmoid\_dual\_iiqc.m} checks the square-factor inverses, the $\mathbf D$ map, weight duality, coercivity, and the identity input-coordinate map. It tests actual PNAS sigmoid increments and arbitrary diagonal slope signals, using held inputs and exact matrix-exponential integration of the quadratic supplies. Every held interval is checked at its midpoint and endpoint. The tested cases include two/four channels, unequal positive pair weights, strictly proper $H$ ($\beta=0$), and the earlier biproper choice $\beta=0.5$. It also tests the stronger lower bound in Proposition~\ref{prop:alternative} and both complementary supplies for the matching plant in the time domain. Numerical tests supplement the proof; the proof establishes the inequality for every finite horizon.

The retained run passed all 20 signal cases and 20,020 finite-prefix checks. The largest transfer duality-map error was $1.4\times10^{-15}$. The largest time-domain complementary identity errors were $3.5\times10^{-18}$ (primal) and $7.2\times10^{-15}$ (dual). Both the strictly proper and biproper cases passed. The script saves \texttt{dual\_verification\_results.mat}; console output is retained in \texttt{dual\_verification\_run.log}.
